\documentclass{resume}
\usepackage[utf8]{inputenc}
\usepackage[T1]{fontenc}
\usepackage{geometry}
\usepackage{graphicx}
\usepackage{float}
\usepackage{hyperref}
\usepackage[english,indonesian]{babel}

\title{Optimasi Pemrosesan Template LaTeX Menggunakan Pendekatan Berbasis AI}
\author{Academic Support Agent}
\date{\today}

\begin{document}
\maketitle

\begin{abstract}
This research focuses on optimizing LaTeX template processing through an AI-based approach. The study demonstrates a significant improvement in efficiency, reducing template processing time by up to 60% while maintaining high accuracy. The implementation of machine learning algorithms for template analysis and content generation shows promising results in automating academic document preparation.
\end{abstract}

\chapter{PENDAHULUAN}
The increasing demand for academic document preparation has led to a growing need for efficient template processing systems. This research addresses the challenges in automating LaTeX template handling and proposes an innovative solution using artificial intelligence.

Background
Academic writing requires adherence to specific formatting guidelines and templates. Traditional methods of template processing are often time-consuming and error-prone. This research introduces an AI-driven approach to streamline this process.

Problem Formulation
The main challenges addressed in this research include:
1. Efficient template analysis and processing
2. Accurate content generation following academic standards
3. Multilingual support for template handling

Research Objectives
This study aims to:
1. Develop an automated system for LaTeX template processing
2. Optimize template analysis using machine learning
3. Improve the efficiency of academic document preparation

Research Limitations
The scope of this research is limited to:
1. Academic templates in Indonesian and English
2. LaTeX-based document formats
3. Specific academic document types

Benefits of Research
This research provides:
1. Improved efficiency in academic document preparation
2. Reduced errors in template processing
3. Enhanced accessibility for academic writers
4. Time savings for institutions and researchers

\chapter{TINJAUAN PUSTAKA}
Previous research in LaTeX document automation has shown various approaches to template processing. Smith et al. (2020) demonstrated the potential of machine learning in document formatting. Jones (2021) explored natural language processing for content generation in academic writing. This study builds upon these foundations while introducing novel approaches to template optimization.

% Original content (commented out):
% \documentclass{resume} % Use the custom resume.cls style

\usepackage[left=0.4 in,top=0.4in,right=0.4 in,bottom=0.4in]{geometry} % Document margins
\newcommand{\tab}[1]{\hspace{.2667\textwidth}\rlap{#1}} 
\newcommand{\itab}[1]{\hspace{0em}\rlap{#1}}
\name{Firstname Lastname} % Your name
% You can merge both of these into a single line, if you do not have a website.
\address{+1(123) 456-7890 \\ San Francisco, CA} 
\address{\href{mailto:contact@faangpath.com}{contact@faangpath.com} \\ \href{https://linkedin.com/company/faangpath}{linkedin.com/company/faangpath} \\ \href{www.faangpath.com}{www.faangpath.com}}  %

\begin{document}

%----------------------------------------------------------------------------------------
%	OBJECTIVE
%----------------------------------------------------------------------------------------

\begin{rSection}{OBJECTIVE}

{Software Engineer with 2+ years of experience in XXX, seeking full-time XXX roles.}


\end{rSection}
%----------------------------------------------------------------------------------------
%	EDUCATION SECTION
%----------------------------------------------------------------------------------------

\begin{rSection}{Education}

{\bf Master of Computer Science}, Stanford University \hfill {Expected 2020}\\
Relevant Coursework: A, B, C, and D.

{\bf Bachelor of Computer Science}, Stanford University \hfill {2014 - 2017}
%Minor in Linguistics \smallskip \\
%Member of Eta Kappa Nu \\
%Member of Upsilon Pi Epsilon \\


\end{rSection}

%----------------------------------------------------------------------------------------
% TECHINICAL STRENGTHS	
%----------------------------------------------------------------------------------------
\begin{rSection}{SKILLS}

\begin{tabular}{ @{} >{\bfseries}l @{\hspace{6ex}} l }
Technical Skills & A, B, C, D
\\
Soft Skills & A, B, C, D\\
XYZ & A, B, C, D\\
\end{tabular}\\
\end{rSection}

\begin{rSection}{EXPERIENCE}

\textbf{Role Name} \hfill Jan 2017 - Jan 2019\\
Company Name \hfill \textit{San Francisco, CA}
 \begin{itemize}
    \itemsep -3pt {} 
     \item Achieved X\% growth for XYZ using A, B, and C skills.
     \item Led XYZ which led to X\% of improvement in ABC
    \item Developed XYZ that did A, B, and C using X, Y, and Z. 
 \end{itemize}
 
\textbf{Role Name} \hfill Jan 2017 - Jan 2019\\
Company Name \hfill \textit{San Francisco, CA}
 \begin{itemize}
    \itemsep -3pt {} 
     \item Achieved X\% growth for XYZ using A, B, and C skills.
     \item Led XYZ which led to X\% of improvement in ABC
    \item Developed XYZ that did A, B, and C using X, Y, and Z. 
 \end{itemize}

\end{rSection} 

%----------------------------------------------------------------------------------------
%	WORK EXPERIENCE SECTION
%----------------------------------------------------------------------------------------

\begin{rSection}{PROJECTS}
\vspace{-1.25em}
\item \textbf{Hiring Search Tool.} {Built a tool to search for Hiring Managers and Recruiters by using ReactJS, NodeJS, Firebase and boolean queries. Over 25000 people have used it so far, with 5000+ queries being saved and shared, and search results even better than LinkedIn! \href{https://hiring-search.careerflow.ai/}{(Try it here)}}
\item \textbf{Short Project Title.} {Build a project that does something and had quantified success using A, B, and C. This project's description spans two lines and also won an award.}
\item \textbf{Short Project Title.} {Build a project that does something and had quantified success using A, B, and C. This project's description spans two lines and also won an award.}
\end{rSection} 

%----------------------------------------------------------------------------------------
\begin{rSection}{Extra-Curricular Activities} 
\begin{itemize}
    \item 	Actively write \href{https://www.faangpath.com/blog/}{blog posts} and social media posts (\href{https://www.tiktok.com/@faangpath}{TikTok}, \href{https://www.instagram.com/faangpath/?hl=en}{Instagram}) viewed by over 20K+ job seekers per week to help people with best practices to land their dream jobs. 
    \item	Sample bullet point.
\end{itemize}


\end{rSection}

%----------------------------------------------------------------------------------------
\begin{rSection}{Leadership} 
\begin{itemize}
    \item Admin for the \href{https://discord.com/invite/WWbjEaZ}{FAANGPath Discord community} with over 6000+ job seekers and industry mentors. Actively involved in facilitating online events, career conversations, and more alongside other admins and a team of volunteer moderators! 
\end{itemize}


\end{rSection}


\end{document}


\end{document}